\section{Reconstruction after invisible endpoint elimination}\label{sec:endpoints}
Use \eqref{eq:intro-profile}. For $I\subseteq[e]=\{1,\ldots,e\}$ put
\begin{align}\label{eq:active-data}
 z_I(x)&=-C_{II}^{-1}B_I^{\mathsf T}x,\quad z_{I^c}(x)=0,\notag\\
 K_I&=\{x\ge0:z_I(x)\ge0,\ B_j^{\mathsf T}x+C_{jI}z_I(x)\ge0\ (j\notin I)\},\\
 S_I&=A-B_IC_{II}^{-1}B_I^{\mathsf T}.\notag
\end{align}
Empty blocks have their natural meanings, in particular $S_\varnothing=A$.
Strict positivity of $C$ makes the endpoint minimizer unique. The stationarity and complementary inequalities show that the cones $K_I$ cover the positive orthant and
$G_Q(x)=x^{\mathsf T}S_Ix$ on $K_I$. The formulas agree on overlaps. Thus $G_Q$ is continuous and is quadratic on finitely many polyhedral regions. We use the quadratic coefficient $S_I$, namely one half the Hessian, as the label of a region.

\begin{definition}[Full visibility]\label{def:visible}
A representation is fully visible if $B$ has column rank $e$ and for every $I\subseteq[e]$ there is an $x\in(0,\infty)^k$ with
\[
 z_I(x)>0,\qquad B_j^{\mathsf T}x+C_{jI}z_I(x)>0\quad(j\notin I).
\]
The first inequality has no entries when $I$ is empty. These strict inequalities give an open region for each active set.
\end{definition}
Let $\mathcal H(G)$ be the set of distinct symmetric matrices $S$ such that $G(x)=x^{\mathsf T}Sx$ on some nonempty open subset of $(0,\infty)^k$. The set, and the open regions carrying its labels, are determined by the function itself.

\begin{lemma}[The visible Hessian lattice]\label{lem:lattice}
For a fully visible representation, $\mathcal H(G_Q)=\{S_I:I\subseteq[e]\}$ consists of exactly $2^e$ matrices, and
\begin{equation}\label{eq:lattice}
 S_I\ge S_J\quad\Longleftrightarrow\quad I\subseteq J.
\end{equation}
For $I\subseteq J$, $\rank(S_I-S_J)=|J\setminus I|$. In particular $A$ is the unique largest local matrix and $S_{[e]}$ the unique smallest one.
\end{lemma}
\begin{proof}
Set $P_I=A-S_I$. A block Schur complement, with $J=I\cup H$, gives
\[
 P_J-P_I=(B_H-B_IC_{II}^{-1}C_{IH})
 (C_{HH}-C_{HI}C_{II}^{-1}C_{IH})^{-1}
 (B_H-B_IC_{II}^{-1}C_{IH})^{\mathsf T}.
\]
The middle matrix is positive definite and the outer matrix has column rank $|H|$, because the columns of $B$ are independent. This proves the forward implication and its rank assertion. Conversely $P_I\le P_J$ implies $\ran P_I\subseteq\ran P_J$: every vector annihilated by $P_J$ has zero quadratic value for $P_I$ and is annihilated by it as well. Since $\ran P_I=\ran B_I$, independence of the columns forces $I\subseteq J$. This proves \eqref{eq:lattice} and distinctness. Full visibility makes all these polynomials occur on open sets. No additional polynomial can occur: any open set meets an interior cell of the finite polyhedral cover, and equality of quadratic polynomials on an open set implies equality of their matrices.
\end{proof}

\begin{theorem}[All-dimensional endpoint reconstruction]\label{thm:reconstruct}
Suppose $G=G_Q$ has a fully visible representation with $e\ge1$ endpoints. The function $G$ determines its minimal endpoint dimension, which is $e$, and every representation with exactly $e$ endpoints as follows.

Read $A=\max\mathcal H(G)$ and $S=\min\mathcal H(G)$. The $e$ matrices covered by $A$ in the ordered set \eqref{eq:lattice} have rank-one deficits
\[
 A-S_{\{j\}}=u_ju_j^{\mathsf T}.
\]
Choose $x_0$ in the open region with local matrix $A$, avoiding all cell boundaries, and orient each $u_j$ by $u_j^{\mathsf T}x_0>0$. Put $U=(u_1,\ldots,u_e)$ and $U^+=(U^{\mathsf T}U)^{-1}U^{\mathsf T}$. Then
\begin{equation}\label{eq:endpoint-reconstruct}
 \overline C=\bigl[U^+(A-S)(U^+)^{\mathsf T}\bigr]^{-1},\qquad
 Q_{\rm norm}=\begin{pmatrix}A&U\\U^{\mathsf T}&\overline C\end{pmatrix},
 \qquad \diag\overline C=\mathbf1.
\end{equation}
These data are unique up to simultaneous permutation of the $u_j$ and endpoint coordinates. A lexicographic ordering of the distinct oriented columns fixes a representative. All $e$-endpoint representations are its positive diagonal endpoint congruences and permutations. Their continuous fibre dimension is $e$.

The equivalent distance representations
\begin{equation}\label{eq:distance-representation}
 G(x)=x^{\mathsf T}Sx+\dist(Tx,K)^2,
 \qquad K=L\R_+^e,
\end{equation}
of minimal dimension have the same $S$ and the same pair $(T,K)$ up to an orthogonal change of ambient coordinates. The normalized Gram matrix of the generating rays is $\overline C$.
\end{theorem}
\begin{proof}
In the original representation set $D_0=\diag(\sqrt{C_{11}},\ldots,\sqrt{C_{ee}})$. The rank-one deficits give $u_j=\pm B_j/\sqrt{C_{jj}}$. On the strict empty-active region all $B_j^{\mathsf T}x$ are positive, so the prescribed orientation selects the positive sign. This orientation does not depend on $x_0$ in that region. Thus $U=BD_0^{-1}$, up to permutation. The normalized coupling is $\overline C=D_0^{-1}CD_0^{-1}$ and
\[
 A-S=U\overline C^{-1}U^{\mathsf T}.
\]
Multiplication by $U^+$ and its transpose gives \eqref{eq:endpoint-reconstruct}. The normalized matrix is positive definite by congruence from $Q$. Its diagonal condition follows directly from the construction.

We next prove minimality without presuming visibility of an alternative representation. In any representation with $e'$ endpoints, the range of every difference of two local quadratic matrices lies in the range of its cross block $B'$, by \eqref{eq:active-data}. The intrinsic difference $A-S$ has rank $e$. Hence $e\le\rank B'\le e'$. No smaller representation is possible.

If $e'=e$, then $B'$ has full column rank. The alternative representation has at most $2^e$ active-set polynomials, while $G$ has exactly $2^e$ distinct local polynomials. Therefore every active set supplies a distinct full-dimensional cell. Each such cell has a strict point: a zero-coordinate or zero-multiplier equality is a nonzero linear equation because the columns of $B'$ are independent, so finitely many such hyperplanes cannot cover an open cell. The alternative representation is fully visible. The preceding reconstruction applies to it and gives exactly the same normalized matrix, up to permutation. Thus every representation is obtained from \eqref{eq:endpoint-reconstruct} by
\[
 Q_D=\begin{pmatrix}A&UD\\DU^{\mathsf T}&D\overline C D\end{pmatrix},
 \qquad D=\diag(d_1,\ldots,d_e),\quad d_j>0,
\]
and a permutation. Conversely changing $z$ by $Dz$ preserves the nonnegative orthant and the infimum, so all these matrices represent $G$. The parameters $d_j$ are freely and injectively specified by the endpoint diagonal. Permutations are discrete. This proves the fibre assertion.

For the distance statement, completing the square gives
$S=A-BC^{-1}B^{\mathsf T}$, $C=L^{\mathsf T}L$ and
$B=-T^{\mathsf T}L$. Conversely every such square expansion gives the positive matrix with $A=S+T^{\mathsf T}T$. Under the preceding congruence the generators become $LD$, up to permutation, which leave the cone unchanged. Any two invertible factors of a fixed Gram matrix differ by a left orthogonal map, and $B=-T^{\mathsf T}L$ forces the same orthogonal map on $T$. Thus $(T,K)$ has exactly the stated ambiguity. Normalizing the lengths of the rays gives $\overline C$.
\end{proof}

\begin{proposition}[An open visible class in every dimension]\label{prop:positive-span}
Suppose $B$ has column rank $e$ and its kernel on the transpose contains a strictly positive vector:
\[
 \text{there is }h\in(0,\infty)^k\text{ with }B^{\mathsf T}h=0.
\]
Then the representation is fully visible, for every $C>0$ and every $A$ with $A-BC^{-1}B^{\mathsf T}>0$. These conditions describe an open nonempty class when $k\ge e+1$. In particular one may take $k=e+1$, let the first $e$ rows of $B$ form $I_e$ and its last row be $-\mathbf1^{\mathsf T}$, choose any $C>0$, and put $A=BC^{-1}B^{\mathsf T}+I_k$.
\end{proposition}
\begin{proof}
Surjectivity of $B^{\mathsf T}$ gives, for any $y\in\R^e$, a solution of $B^{\mathsf T}x=y$. Adding a sufficiently large multiple of $h$ makes $x$ strictly positive. Thus $B^{\mathsf T}(0,\infty)^k=\R^e$.
For an active set $I$, choose $z_I>0$, $z_{I^c}=0$, $w_I=0$, and $w_{I^c}>0$. There exists an $x>0$ with $B^{\mathsf T}x=-Cz+w$. These are precisely the strict stationarity and complementary inequalities, so $I$ is visible.

The property persists under small perturbations of $B$. Indeed, solve the small residual in $B'^{\mathsf T}h$ using a bounded right inverse of $B'^{\mathsf T}$; the correcting vector tends to zero, so the corrected kernel vector remains positive. Rank, positive definiteness, and the Schur inequality are open conditions. The displayed matrix has positive kernel vector $\mathbf1$, and proves nonemptiness. Additional rows for $k>e+1$ may initially be zero and the same positive-kernel argument applies. Thus this is an open family of genuinely coupled endpoint problems, not merely the diagonal case.
\end{proof}

\begin{proposition}[Finite redundancy criterion]\label{prop:redundancy}
For an arbitrary positive definite $Q$, deleting endpoint coordinate $j$ preserves $G_Q$ if and only if, for every active cone $K_I^{(-j)}$ of the reduced endpoint problem,
\begin{equation}\label{eq:redundancy}
 \bigl(B_j^{\mathsf T}-C_{jI}C_{II}^{-1}B_I^{\mathsf T}\bigr)x\ge0
 \quad\hbox{for every }x\in K_I^{(-j)}.
\end{equation}
Each condition is membership of a specified linear functional in the dual of a polyhedral cone and can be tested by a finite linear feasibility problem. Independently of visibility, the invariant
\[
 \delta(G)=\dim\operatorname{span}\{\ran(S-S'):S,S'\in\mathcal H(G)\}
\]
is a lower bound for every endpoint dimension. In the fully visible class it equals the minimal dimension $e$.
\end{proposition}
\begin{proof}
The reduced minimizer is feasible for the full problem with its $j$th coordinate zero. Its gradient in that coordinate, divided by two, is the left side of \eqref{eq:redundancy}. The reduced stationarity conditions already hold in all other coordinates. The additional nonnegative gradient is therefore necessary and sufficient for full optimality. Strict convexity shows that equality of the two value functions for every $x$ is equivalent to this condition at the unique reduced minimizer for every $x$. Substituting the finitely many active formulas gives \eqref{eq:redundancy}, including boundaries. The range bound follows because every local matrix is $A-B_IC_{II}^{-1}B_I^{\mathsf T}$; every difference has range in $\ran B$. In the visible class $A-S_{[e]}$ already has range of dimension $e$, giving equality.
\end{proof}
The coordinate-deletion criterion does not assert that repeated coordinate deletion finds a minimal representation in every nonvisible problem. Theorem~\ref{thm:reconstruct} proves global minimality and the complete minimal-dimensional fibre on its stated open class. Outside that class the finite equality test by intersections of active cones remains valid, but a full classification of all nonminimal representations is a separate question.

